% CS 519 Project Status Update -- April 7, 2026
% Compile with: pdflatex status_update.tex
% Requires neurips_2024.sty in the same directory.
\documentclass{article}
\usepackage[preprint]{neurips_2024}

\usepackage{amsmath}
\usepackage{booktabs}
\usepackage{graphicx}
\usepackage{hyperref}
\usepackage{microtype}

\title{Inductive Graph-Based Skill Recommendation\\
       using Workforce Knowledge Graphs\\[4pt]
       \large CS 519 Project Status Update --- April 7, 2026}

\author{
  Akshaj Kurra Satishkumar \quad Monu Kumar \\
  Department of Computer Science, University of Illinois Chicago \\
  \texttt{akurr@uic.edu} \quad \texttt{mkuma47@uic.edu}
}

\begin{document}
\maketitle

%-----------------------------------------------------------------------
\section{Introduction and Motivation}
%-----------------------------------------------------------------------

Modern labor markets evolve faster than the models trained to serve them.
New occupational roles emerge without historical training data, yet
practitioners need skill guidance from day one.
Classical collaborative-filtering methods---cosine similarity on skill
vectors, Jaccard overlap on binary skill sets---represent each occupation
entirely from its observed training edges.
When an occupation has no training edges (the \emph{cold-start} setting),
these methods produce all-zero score vectors and cannot recommend anything.
Node2Vec and other transductive embedding methods face the same limitation:
all node embeddings are learned as lookup tables at training time and cannot
be updated for new nodes at inference.

We address this by framing occupation--skill recommendation as weighted
link prediction on a heterogeneous knowledge graph and applying an
\emph{inductive} graph neural network (BipartiteSAGE).
Because GraphSAGE learns a parameterized aggregation function rather than a
fixed embedding table, it can produce embeddings for previously unseen
occupations by aggregating over node features alone---enabling cold-start
generalization that is impossible for the transductive baselines.

%-----------------------------------------------------------------------
\section{Problem Formulation}
%-----------------------------------------------------------------------

Let $\mathcal{O}$ be the set of occupations and $\mathcal{S}$ the set of
workforce features (skills, abilities, work activities, knowledge).
We construct a weighted bipartite graph
$G = (\mathcal{O} \cup \mathcal{S},\,\mathcal{E})$
where each edge $(o,s)\in\mathcal{E}$ carries weight
\[
  w(o,s) = \tfrac{1}{2}\!\left(\tfrac{\mathrm{IM}-1}{4}
           + \tfrac{\mathrm{LV}}{7}\right) \in [0,1],
\]
normalizing O*NET's 1--5 Importance and 0--7 Level ratings.
We retain only edges with $w \geq 0.5$ (high importance/level), yielding
44,619 edges over 894 occupations and 160 feature nodes.

\paragraph{Task 1 -- Skill Recommendation.}
Given the training subgraph, rank all candidate features by relevance for
each occupation.
Evaluated with Recall@$K$ ($K\in\{5,10\}$), NDCG@$K$,
Precision@10, and AUC (one random negative per positive edge, sampled
outside the training set).

\paragraph{Task 2 -- Occupation Similarity.}
Given learned occupation embeddings, retrieve the $k$ nearest occupations
by cosine similarity and measure the fraction sharing the same SOC major
group (first two digits of the O*NET code, e.g., ``11'' for Management).
This evaluates structural coherence of the embedding space without
requiring held-out similarity annotations.

\paragraph{Evaluation splits.}
\textbf{Standard:} 20\% of edges per occupation held out; all occupations
appear in training.
\textbf{Cold-start:} 10\% of occupations ($\approx 89$) fully withheld from
training and evaluated inductively at test time.
Training edges are masked from ranking to prevent trivially re-recommending
observed links.

%-----------------------------------------------------------------------
\section{Methods}
%-----------------------------------------------------------------------

\paragraph{Graph construction.}
We aggregate four O*NET Release 30.1 dimensions---Skills, Abilities, Work
Activities, and Knowledge---prefixing each element name by dimension to keep
features distinct (e.g., \texttt{Skills/Active Listening}).
After weight thresholding, the PyTorch Geometric \texttt{HeteroData} graph
has edge types \texttt{(occupation, requires, feature)} and its reverse.
Occupation features are $|\mathcal{S}|$-dimensional weighted skill profiles
built from training edges only; feature-node features are identity vectors.

\paragraph{Cosine and Jaccard baselines.}
Cosine represents each occupation as a weighted feature vector and predicts
via neighbor-weighted averaging (user-based CF).
Jaccard uses binarized feature sets with intersection-over-union similarity.
Both are transductive; cold-start occupations have all-zero feature vectors
and receive identical, meaningless scores.

\paragraph{Node2Vec~\cite{grover2016node2vec}.}
We project the bipartite graph to a homogeneous space (feature indices
offset by $|\mathcal{O}|$) and train skip-gram embeddings
(walk length 20, context 10, embedding dim 64, 50 epochs).
Scores are $\mathbf{E}_\mathcal{O}\mathbf{E}_\mathcal{S}^\top$.
Transductive by construction; cannot score cold-start occupations.

\paragraph{BipartiteSAGE~\cite{hamilton2017inductive}.}
Two-layer \texttt{HeteroConv} with \texttt{SAGEConv} on both edge
directions, hidden 128, output 64, dot-product decoder.
For cold-start, occupation nodes receive TF-IDF vectors (500 dimensions,
fitted on training occupations only) of their O*NET task description text,
concatenated with the skill-profile features.
Trained with binary cross-entropy loss and one random negative per positive
edge (Adam, lr\,=\,0.01, 100 epochs); the checkpoint with the best Recall@10
on the validation set is retained.

%-----------------------------------------------------------------------
\section{Results and Discussion}
%-----------------------------------------------------------------------

Table~\ref{tab:task1} reports the full Task~1 comparison.
All experiments are complete and reproducible from \texttt{python main.py}.

\begin{table}[h]
\centering
\caption{Task 1 -- Skill Recommendation results after edge-weight thresholding
  ($w\!\geq\!0.5$, 44{,}619 edges).
  Higher is better for all metrics.
  Transductive baselines cannot operate in the cold-start setting.}
\label{tab:task1}
\small
\begin{tabular}{lcccccc}
\toprule
Method & R@5 & R@10 & NDCG@5 & NDCG@10 & P@10 & AUC \\
\midrule
Cosine (std)       & 0.337 & 0.532 & 0.719 & 0.664 & 0.546 & 0.900 \\
Jaccard (std)      & 0.345 & 0.547 & 0.734 & 0.680 & 0.559 & 0.900 \\
Node2Vec (std)     & 0.125 & 0.227 & 0.301 & 0.288 & 0.266 & 0.731 \\
\textbf{GraphSAGE (std)}  & \textbf{0.365} & \textbf{0.575} & \textbf{0.753} & \textbf{0.702} & \textbf{0.576} & \textbf{0.919} \\
\midrule
\textbf{GraphSAGE (cold)} & 0.111 & 0.216 & \textbf{0.934} & \textbf{0.921} & \textbf{0.906} & 0.904 \\
Cosine / Jaccard (cold) & \multicolumn{6}{c}{\textit{not applicable -- cannot score unseen occupations}} \\
\bottomrule
\end{tabular}
\end{table}

\paragraph{Standard split.}
After thresholding to the 44,619 high-salience edges, the evaluation is
meaningfully discriminative.
GraphSAGE achieves the best performance across all metrics
(Recall@10 0.575, NDCG@10 0.702, AUC 0.919), outperforming both cosine and
Jaccard baselines.
Node2Vec performs significantly worse (AUC 0.731), confirming that
random-walk topology embedding---without node features---is insufficient on a
dense bipartite graph where walk patterns are uninformative.

\paragraph{Cold-start split.}
The key finding is the split between ranking \emph{quality} and
\emph{coverage}.
NDCG@5 reaches 0.934, meaning that when GraphSAGE recommends features for
an unseen occupation, the top-5 items are nearly always highly relevant.
However, Recall@5 is 0.111: with $\sim\!50$ gold features per occupation on
average, the top-5 list covers only a small fraction.
This gap is \emph{expected}: the model cannot enumerate all features of an
unseen role from text features alone, but the features it does surface are
correct.
AUC of 0.904 under cold-start confirms a learned global ranking function
generalizes inductively.

\paragraph{Task 2.}
GraphSAGE embeddings achieve FamilyPrecision@5 = 0.264 and
FamilyPrecision@10 = 0.246, meaning roughly 25\% of nearest neighbors in
embedding space share the same SOC major group.
Given that there are 23 SOC families among 894 occupations, random chance
would yield $\approx 4\%$; the learned representation clusters meaningfully
by occupational family.

\paragraph{Graph density and thresholding.}
Prior to thresholding, the graph had 135,366 edges (94\% feature coverage
per occupation), making Recall@K and Precision@K trivially near 1.
The $w\!\geq\!0.5$ threshold retains only high-salience associations (31\%
coverage), creating a realistic recommendation task where sparsity matters.
This threshold is a design choice we will analyze in ablations.

%-----------------------------------------------------------------------
\section{Future Plans}
%-----------------------------------------------------------------------

\begin{itemize}
  \item \textbf{Weighted message passing.}
        Incorporate edge weight $w(o,s)$ directly into the SAGEConv
        aggregation so high-importance links contribute proportionally more.
        This is the clearest structural improvement to the model.

  \item \textbf{Cold-start coverage improvement.}
        Recall@K for cold-start is limited by the number of gold edges
        that can fit in top-$K$.
        We will experiment with larger $K$ (20, 50) and explore whether
        aggregating over semantically similar \emph{training} occupations
        (via TF-IDF nearest neighbors) can increase effective neighborhood
        size at inference time.

  \item \textbf{Hyperparameter ablations and GAT variant.}
        Sweep depth (1--3 layers), hidden dimension (64/128/256), and
        aggregation type (mean vs.\ LSTM).
        Add a graph attention (GAT) variant to assess whether dynamic
        edge weighting further improves over uniform SAGEConv.

  \item \textbf{Threshold sensitivity.}
        Evaluate results under multiple thresholds ($w \geq 0.3, 0.4, 0.5$)
        and report how graph density affects each method's performance.
        This contextualizes the comparison and addresses the density issue
        systematically.
\end{itemize}

%-----------------------------------------------------------------------
\begin{thebibliography}{9}
\bibitem{hamilton2017inductive}
W.~L. Hamilton, Z.~Ying, and J.~Leskovec.
Inductive representation learning on large graphs.
\textit{NeurIPS}, 2017.

\bibitem{grover2016node2vec}
A.~Grover and J.~Leskovec.
node2vec: Scalable feature learning for networks.
\textit{ACM SIGKDD}, 2016.

\bibitem{onet301}
National Center for O*NET Development.
O*NET Database Release 30.1.
U.S.\ Department of Labor, 2024.
\texttt{https://www.onetcenter.org}
\end{thebibliography}

\end{document}
